% Methodology: Counterfactuals
% Technical annex to Watling, S. and Breach, A. (2023), "The housebuilding
% crisis: The UK's 4 million missing homes", Centre for Cities.
%
% Corrected edition, September 2026. The February 2023 annex is revised so that
% every equation and worked figure matches the code that reproduces the report's
% results: code/Counterfactual Replication.R for the counterfactuals and
% code/Summary Replication.R for the housing stock estimate. Worked figures were
% recomputed from data/Combined.csv, and the Stage 7 formulas reproduce all 13
% rows of the code's output. The final section lists the changes.
%
% Build: latexmk -pdf methodology.tex

\documentclass[11pt,a4paper]{article}

\usepackage[T1]{fontenc}
\usepackage[utf8]{inputenc}
\usepackage{lmodern}
\usepackage{microtype}
\usepackage[margin=2.5cm]{geometry}
\usepackage{amsmath,amssymb}
\usepackage{booktabs}
\usepackage{enumitem}
\usepackage[hidelinks]{hyperref}

\hypersetup{
  pdftitle  = {Methodology: Counterfactuals},
  pdfauthor = {Samuel Watling and Anthony Breach}
}

\setlength{\parindent}{0pt}
\setlength{\parskip}{0.6em}

% ---- Notation -----------------------------------------------------------------
\newcommand{\priv}{\mathrm{Priv}}
\newcommand{\pub}{\mathrm{Pub}}
\newcommand{\PC}{\mathit{PC}}
\newcommand{\QK}{\mathit{QK}}
\newcommand{\code}[1]{\texttt{#1}}

\title{Methodology: Counterfactuals\\[0.5em]
  \large\emph{The housebuilding crisis: The UK's 4 million missing homes}}
\author{Samuel Watling and Anthony Breach\\[0.3em]\normalsize Centre for Cities}
\date{February 2023\\[0.2em]\small Corrected edition, September 2026}

\begin{document}

\maketitle

\begin{center}
\begin{minipage}{0.9\textwidth}\small
Technical annex to the report, originally published by Centre for Cities,
\copyright\ Centre for Cities 2023. This corrected edition revises the equations and
worked figures so that they match the code that reproduces the report's results:
\code{code/Counterfactual Replication.R} for the counterfactuals and
\code{code/Summary Replication.R} for the housing stock estimate. Every worked figure
has been recomputed from \code{data/Combined.csv}. The final section lists the changes
from the February 2023 text.
\end{minipage}
\end{center}

% ==============================================================================
\section*{Introduction}

In the report \textbf{The housebuilding crisis: The UK's 4 million missing homes},
housing outcomes are defined as the number of homes per (thousand) people. This is a
simple metric for general comparisons in housing quality and affordability.

One of the goals of the report was to calculate exactly how much underbuilding had
taken place in the UK since 1955 compared to other countries, and how this affected
outcomes in the UK.

Comparisons that simply applied the housebuilding rates of other countries to the UK
generated unbuilt backlogs that were too large, as other countries in 1955 had worse
outcomes and higher population growth, and therefore needed to build more.

More reasonable estimates could be produced by calculating the housebuilding required
to reach the number of homes per person that other European countries had in 2015, but
these were still unsatisfactory, for two reasons.

First, investigating differences in housebuilding by tenure was not possible with this
method, despite being a key aspect of the report.

And second, inconsistencies in the data between countries make it difficult to draw
comparisons based on the number of dwellings recorded in 2015. Both the UK, and wherever
possible the UN data, record the total housing stock until 2000. However, the
statistical agencies of many other countries today only record primary residences, and
omit second homes and vacant stock. Therefore, using the publicly available 2015
dwellings per person ratios from national statistics agencies to estimate UK
underbuilding from 1955 generates underestimates and inconsistent results between
countries.

To solve this, Centre for Cities used a new method that combined both approaches. The
report's estimates of the total amount of underbuilding in the UK from 1955 to 2015
relative to other countries are the difference between the number of homes the UK would
have had if it had built at the rate of another country and the number it actually had,
controlling for the difference in the initial stock of homes per person and for
different population growth rates.

The number of homes per person in any year is
\[
  \text{Homes per person}_t
  = \frac{\text{Previous year's housing stock} + \text{Net additions}_t}{\text{Population}_t},
\]
with
\[
  \text{Net additions}_t = \text{Homes built}_t - \text{Homes demolished}_t.
\]

This calculation is done twice: first to estimate the actual British net additions every
year, and then again for each ``counterfactual'' version of the UK, adjusted for each
country comparison (e.g.\ the Netherlands). The difference between the final values for
the actual UK and each counterfactual is the estimate of the housebuilding that would have
taken place in the UK had it followed the housing policy approach of that country.

Part~I explains the calculation for the actual UK, including how many homes the private
and public sectors added from 1955 to 2015. Part~II explains how it is adjusted for each
counterfactual.

% ==============================================================================
\section*{Data and notation}

Years are indexed $t = 1, \dots, T$ for 1955 to 2015, so $T = 61$. Worked figures are in
thousands of homes unless stated. They are rounded for display only: the code rounds
nothing until its final output, so displayed figures may not add exactly. For each
country the data provide:
\begin{itemize}
  \item $S_t$, the estimated housing stock for year $t$ (\code{X5 (Estimate)}),
        constructed as described under \emph{Interpolation};
  \item $B_t$, gross completions (\code{Y7});
  \item $\theta_t^{\priv}$ and $\theta_t^{\pub}$, the private and public shares of
        completions (columns \code{4} and \code{1});
  \item population in millions, written $P_t$ for the UK and $Q_t$ for a comparator
        country, where $P_0$ and $Q_0$ are the 1954 values;
  \item homes per 1{,}000 people in 1955, written $\PC_1$ for the UK and $\QK_1$ for a
        comparator.
\end{itemize}
Building rates are shares of the estimated stock, by tenure:
\[
  b_t^{\priv} = \frac{\theta_t^{\priv} B_t}{S_t}, \qquad
  b_t^{\pub}  = \frac{\theta_t^{\pub} B_t}{S_t}, \qquad
  b_t = b_t^{\priv} + b_t^{\pub},
\]
for the UK, and $c_t^{\priv}$, $c_t^{\pub}$ and $c_t$ are defined in the same way for a
comparator. The two tenure shares are used as recorded and do not always sum to one: they
fall to as little as 0.37 for Finland in 1955--79, and range from 0.84 to 1.05 for
Switzerland in 1955--78. In those years the building counted by the method differs from
recorded completions.

The UK demolition rate is $d_t = E_t / S_t$, where $E_t$ is reported demolitions. In the
10 years without a report, $E_t$ is completions multiplied by the ratio of demolitions to
completions, interpolated across gaps of up to two years and otherwise carried forward
from the most recent reported year.

% ==============================================================================
\section*{Part I -- Calculating the actual British values}

Four steps are required:
\begin{enumerate}
  \item The housing stock for each year from 1955 to 2015.
  \item The net change in housing stock between 1955 and 2015.
  \item Gross housebuilding and gross demolitions from 1955 to 2015.
  \item Net additions by tenure from 1955 to 2015.
\end{enumerate}

The UK is put through exactly the same calculation as each counterfactual, so that the
two are consistent. The resulting stock is a model value rather than the national
statistic: for 2015 it is 27.65 million against an estimated 28.28 million, 2.2 per cent
lower.

\subsection*{Stage 1: The housing stock for each year from 1955 to 2015}

The calculation starts from the estimated stock in 1955, $H_0 = S_1 = 15{,}418$. Each
year the stock grows by that year's building rate net of its demolition rate:
\[
  H_t = H_{t-1}(1 + b_t - d_t),
  \qquad\text{so}\qquad
  H_t = H_0 \prod_{s=1}^{t} (1 + b_s - d_s).
\]
In 1955 and 1956 the rates are $b_1 = 2.13\%$, $d_1 = 0.36\%$, $b_2 = 1.98\%$ and
$d_2 = 0.41\%$, which give $H_1 = 15.69$ million and $H_2 = 15.94$ million. Continuing to
2015 gives $H_T = 27.65$ million, and dividing by population gives the number of homes per
person in every year.

\subsection*{Stage 2: The net change in housing stock between 1955 and 2015}

The net change is the difference between the final and initial stock, which is the sum
of each year's net additions:
\[
  \Delta H = H_T - H_0 = \sum_{t=1}^{T} H_{t-1}(b_t - d_t).
\]
For the UK, $\Delta H = 27{,}651.0 - 15{,}418.0 = 12{,}233.0$.

\subsection*{Stage 3: Gross housebuilding and gross demolitions}

The net change in total dwellings is not enough to estimate additions by tenure: gross
housebuilding and gross demolitions are needed. The code measures both against the stock
at the \emph{end} of each year, $H_t$, and splits each year's demolitions by that year's
tenure shares:
\[
  G^{\priv} = \sum_{t=1}^{T} H_t\, b_t^{\priv}, \qquad
  E^{\priv} = \sum_{t=1}^{T} H_t\, \theta_t^{\priv} d_t, \qquad
  N^{\priv} = G^{\priv} - E^{\priv},
\]
and likewise for public building, with totals $G = G^{\priv} + G^{\pub}$,
$E = E^{\priv} + E^{\pub}$ and $N = N^{\priv} + N^{\pub}$. For the UK,
\[
  N = G - E = 15{,}512.1 - 3{,}151.2 = 12{,}360.9.
\]

Because each year's rates are applied to the end-of-year stock rather than the
start-of-year stock that generated them, $N$ exceeds the net change in the stock by
\[
  \varepsilon = N - \Delta H = 12{,}360.9 - 12{,}233.0 = 127.9.
\]
When the tenure shares sum to one, this is exactly
$\varepsilon = \sum_{t} H_{t-1}(b_t - d_t)^2$, a second-order term.

\subsection*{Stage 4: Net additions by tenure}

Splitting net change in stock by tenure requires further assumptions about demolitions
(and, for the counterfactuals, population growth). Demolitions are allocated each year to
\textbf{the same tenure proportion as the homes built in that year}; in Part~II, the
population and initial-stock adjustments are allocated by the comparator's tenure mix of
gross building.

British data is available from 1955 to 2015 for gross housebuilding by tenure and gross
demolitions, but not the necessary data to measure demolitions by tenure. Although for the
latter part of the period some data does exist on demolitions by tenure of final occupier
for England, data does not exist on demolitions by tenure of \emph{construction}. As this
report is concerned with the contribution of the private and public sector to total net
additions of stock, it is the latter we are interested in.

To take an example -- consider a local authority which built ten council houses, and four
of them were sold to tenants under Right to Buy. One remaining council tenancy property
and a Right to Buy property were both then demolished. The net change in tenure by
occupation is an additional five council dwellings and three homeowner dwellings. But when
considering the net change by construction, the public sector added eight dwellings,
despite Right to Buy and the demolition of an owner-occupied dwelling.

As data on demolitions by tenure is tracked by final occupier, not by construction, this
presents a problem for the methodology. Former council homes sold under Right to Buy and
former private properties sold to Registered Providers will have their demolitions
recorded incorrectly. The report must make an assumption about demolitions for the entire
period.

Demolitions were higher in the postwar period (specifically the 1960s) and concentrated
on private sector stock in slum clearances. Demolitions since 1980 were lower over a
longer time period, and primarily were of social housing and private properties often
built by the public sector (e.g.\ leaseholders on council estate regenerations).

Allocating each year's demolitions to that year's tenure mix puts 58.6 per cent of UK
demolitions from 1955 to 2015 on private-built homes. This is below the private share of
gross building (64.4 per cent) because demolition was heaviest, relative to building, in
years when public building made up a larger share. The individual results in each year
are skewed, but the cumulative results are satisfactory.

Finally, the tenure totals are scaled so that they add up to the net change in the stock,
removing $\varepsilon$ in proportion to each tenure's net building:
\[
  A^{\priv} = N^{\priv}\,\frac{\Delta H}{N}, \qquad
  A^{\pub}  = N^{\pub}\,\frac{\Delta H}{N}, \qquad
  A^{\priv} + A^{\pub} = \Delta H.
\]

\begin{center}
\begin{tabular}{lrrr}
  \toprule
  UK, 1955--2015 (thousands) & Private & Public & Total\\
  \midrule
  Gross building, $G$  & 9,985.5 & 5,526.7 & 15,512.1\\
  Demolitions, $E$     & 1,847.7 & 1,303.5 &  3,151.2\\
  Net building, $N$    & 8,137.7 & 4,223.2 & 12,360.9\\
  Net additions, $A$   & 8,053.5 & 4,179.5 & 12,233.0\\
  \bottomrule
\end{tabular}
\end{center}

The UK therefore added 8.05 million private and 4.18 million public homes from 1955 to
2015, a split of 66:34.

% ==============================================================================
\section*{Part II -- Calculating counterfactual values\\ for British housebuilding}

Having calculated the required figures for Britain, we now move to the method used to
make credible adjustments for other countries.

The central estimate we are making is \textbf{how many extra homes the UK would have had
in 2015 if it had built at the rate of another Western European country from 1955.} As
accurate values for the net change in housing stocks both in general and by tenure are
not available, this is calculated using each country's gross building and UK demolitions.

However, simply repeating the British method with another country's building rates will
not produce a credible result. This is because:

\textbf{Other countries have different population growth rates compared to Britain.} If
the population growth rate is different to the UK's, then the same gross increase will not
produce the same change in the number of homes per person. As all countries except
Austria and Belgium had higher population growth than the UK, adjusting for this makes the
average estimate more conservative.

\textbf{Other countries began the period with a different per person stock of homes
compared to Britain.} The average western European country began the period with a
smaller stock of homes per person than the UK, which means some of their housebuilding
after 1955 will have been to catch up to a housing stock that the UK had already built
before 1955. Adjusting for this makes the estimate more conservative.

Therefore, these effects need to be controlled for when calculating the counterfactual
value of the United Kingdom, so that it has:
\begin{itemize}
  \item the same initial housing stock to population ratio as the counterfactual country;
  \item the same population growth rate as the counterfactual country.
\end{itemize}

\subsection*{Methodology}

A counterfactual UK that followed a comparator country (e.g.\ the Netherlands) is built in
the same way as the actual UK in Part~I, with three adjustments:
\begin{enumerate}
  \item its initial stock gives it the comparator's 1955 ratio of homes per person;
  \item its stock is rescaled each year by the difference between UK and comparator
        population growth;
  \item it builds at the comparator's rates, and demolishes at the UK's rate adjusted for
        how many more or fewer homes per person it has than the actual UK.
\end{enumerate}
The counterfactual's annual stock is then used for the same tenure calculation as in
Part~I:
\begin{enumerate}[resume]
  \item gross building, demolitions and net building by tenure;
  \item the discrepancy that the population rescaling creates between net building and
        the change in the stock.
\end{enumerate}
Finally:
\begin{enumerate}[resume]
  \item the difference between the counterfactual and the actual UK is split by tenure.
\end{enumerate}

\subsection*{Stage 1: Adjust the counterfactual's initial stock of homes per person}

The counterfactual starts with the UK's 1955 population housed at the comparator's 1955
ratio of homes per person:
\[
  \sigma_0 = P_1\,\QK_1 = H_0\,\frac{\QK_1}{\PC_1},
\]
where the second form uses $\PC_1 = H_0 / P_1$. In 1955 the Netherlands had 233.69
dwellings per 1,000 people against the UK's 302.49, and the UK's population was 50.97
million. So the initial stock of the counterfactual UK based upon the Netherlands is
\[
  \sigma_{\text{Neth}} = 50.97 \times 233.693 = 11{,}911.3.
\]

\subsection*{Stage 2: Net change in stock for the counterfactual UK}

The counterfactual builds at the comparator's rate $c_t$ and demolishes at a rate $D_t$,
which is set out in Stage~4. Before any adjustment for population, its stock would grow
each year by the factor $1 + c_t - D_t$: in the first year, from $\sigma_0$ to
$\sigma_0(1 + c_1 - D_1)$.

\subsection*{Stage 3: Adjust net change in stock for population growth}

Population growth differs between countries every year. This means that an identical
increase in the overall housing stock will not result in an identical increase in the
stock of homes per person. Places with lower population growth do not need to build as
many homes to achieve the same improvements to housing outcomes. Therefore, the change in
population relative to Britain is applied to the housing stock every year.

The code measures each year's population growth relative to that year's population:
\[
  \Delta P_t = 1 + \frac{P_t - P_{t-1}}{P_t}, \qquad
  \Delta Q_t = 1 + \frac{Q_t - Q_{t-1}}{Q_t},
\]
a close approximation to the growth factors $P_t/P_{t-1}$ and $Q_t/Q_{t-1}$. Writing
$R_t = \Delta P_t / \Delta Q_t$, the counterfactual stock $\nu_t$ is
\[
  \nu_t = \nu_{t-1}(1 + c_t - D_t)\,R_t, \quad \nu_0 = \sigma_0,
  \qquad\text{so}\qquad
  \nu_t = \sigma_0 \prod_{s=1}^{t} (1 + c_s - D_s)\,R_s.
\]
Multiplying by $R_t$ means that the counterfactual's homes per person grow as they would
with the comparator's building and population growth, while its stock is always expressed
for the UK's population. The Netherlands grew faster than the UK: in 1955,
$R_1 = 0.9917$, or about $1/1.008$, so the counterfactual stock is scaled down slightly.
For the UK itself, $R_t = 1$ and $D_t = d_t$, so $\nu_t = H_t$.

\subsection*{Stage 4: Adjust demolition rates for the counterfactual's stock of homes}

Accurate and consistent demolition data does not exist for European countries apart from
Switzerland. For all other countries, the demolition rate is set through a two-step
process.

Initially, given that this is a model of Britain's housing, the demolition rates are
assumed to be identical to Britain's, as it is mostly pre-1947 British stock which is
being demolished. This gives a first-round stock $\xi_t$:
\[
  \xi_t = \xi_{t-1}(1 + c_t - d_t)\,R_t, \qquad \xi_0 = \sigma_0.
\]

However, the UK's demolition rate needs to be adjusted for two reasons. First, the UK's
demolitions are concentrated in the 1960s due to a policy choice, with little demolition
earlier or after that. Applying this to the counterfactual version of the UK building at
another rate will have an inconsistent effect on the net change in dwellings. Second, as
the main target of demolitions was unfit pre-modern housing, it is reasonable to assume
that demolition rates would be in proportion to the stock of housing. If the UK had had
less housing than it actually did in 1955, then it likely would have demolished fewer
dwellings in the years after.

Because $\xi_t$ and $H_t$ are both stocks for the UK's population, $\xi_t / H_t$ is the
ratio of the counterfactual's homes per person to the actual UK's. The demolition rate is
moved by one hundredth of the proportional gap:
\[
  D_t = d_t + \frac{1}{100}\left(\frac{\xi_t}{H_t} - 1\right).
\]
A counterfactual with 1 per cent more homes per person than the actual UK demolishes at a
rate 0.01 percentage points higher; one with 1 per cent fewer, at a rate 0.01 percentage
points lower. In 1955 the Netherlands counterfactual has $\xi_1 / H_1 = 0.768$, 23.2 per
cent fewer homes per person than the UK, so its demolition rate is
\[
  D_1 = 0.357\% - 0.232\% = 0.125\%.
\]

The effect of this over time smooths out demolitions from 1955 to 2015. In 1955 the actual
UK has more homes per person than six of the eleven comparators, and by 2015 every
counterfactual has more homes than the actual UK. So most counterfactual versions of the UK
have fewer demolitions at the beginning of the period (and in the 1960s), when they have
fewer homes per person than the UK, and more demolitions towards the end of the period,
when they have overtaken the UK and actual UK demolitions have dropped to very low levels.

Nothing stops $D_t$ falling below zero. It does so for the Ireland counterfactual in 18
years (1978 to 1996), where the negative demolitions add about 239,000 homes, and for the
Netherlands in 1992, where the effect is negligible.

Switzerland is the exception: its demolition data are of high enough quality to use
directly. Its $D_t$ is its own demolitions divided by its own estimated stock, with gaps of
up to two years interpolated and longer gaps filled using the ratio of demolitions to
completions from the most recent reported year (or the first, for years before any
report).

The net result is conservative: across the eleven comparators, it reduces the average
number of additional homes from 4.87 million, with $D_t = d_t$, to 4.25 million.

\subsection*{Stage 5: Gross building, demolitions and net building by tenure}

Part~I is now applied to the counterfactual stock. As there, building and demolitions in
year $t$ are measured against the end-of-year stock $\nu_t$, and demolitions are split by
the comparator's tenure shares:
\[
  J_B^{\priv} = \sum_{t=1}^{T} \nu_t\, c_t^{\priv}, \qquad
  J_D^{\priv} = \sum_{t=1}^{T} \nu_t\, \theta_t^{\priv} D_t, \qquad
  J_N^{\priv} = J_B^{\priv} - J_D^{\priv},
\]
and likewise for public building, with totals $J_B$, $J_D$ and $J_N$. For the
Netherlands counterfactual:

\begin{center}
\begin{tabular}{lrrr}
  \toprule
  Netherlands counterfactual (thousands) & Private & Public & Total\\
  \midrule
  Gross building, $J_B$  & 14,201.1 & 9,881.8 & 24,083.0\\
  Demolitions, $J_D$     &    952.5 &   769.6 &  1,722.2\\
  Net building, $J_N$    & 13,248.6 & 9,112.2 & 22,360.8\\
  \bottomrule
\end{tabular}
\end{center}

\subsection*{Stage 6: Adjust for the population discrepancy}

Net building should equal the change in the stock, but the population rescaling in
Stage~3 changes the stock without any building or demolition. The discrepancy is
\[
  \Gamma = J_N - (\nu_T - \sigma_0) = \sum_{t=1}^{T} \gamma_t,
  \qquad
  \gamma_t = \nu_t\bigl(c_t - (\theta_t^{\priv} + \theta_t^{\pub}) D_t\bigr) - (\nu_t - \nu_{t-1}).
\]
When the tenure shares sum to one,
\[
  \gamma_t = \nu_{t-1}\bigl[1 - R_t\bigl(1 - (c_t - D_t)^2\bigr)\bigr],
\]
which combines the population rescaling with the same end-of-year timing term as
$\varepsilon$ in Part~I. For the UK, $R_t = 1$ and $\Gamma = \varepsilon$. Due to
compounding, the discrepancy runs into millions of homes in the counterfactuals from 1955
to 2015.

For the Netherlands counterfactual, $\nu_T = 30{,}486.7$, so
\[
  \nu_T - \sigma_0 = 30{,}486.7 - 11{,}911.3 = 18{,}575.4,
  \qquad
  \Gamma = 22{,}360.8 - 18{,}575.4 = 3{,}785.4.
\]

\subsection*{Stage 7: Compare to the UK, by tenure}

The aim of the counterfactuals is \textbf{how many extra homes Britain would have of each
tenure had it built at the rate of another European country.} The total is the difference
in final stocks, which already reflects the initial difference in homes per person,
population growth and demolitions:
\[
  \text{Total surplus} = \nu_T - H_T.
\]
For the Netherlands, $\text{Total surplus} = 30{,}486.7 - 27{,}651.0 = 2{,}835.7$.

Part of that difference is a difference in net building, $J_N - N$. The rest is
\[
  \Delta = (\nu_T - H_T) - (J_N - N) = (\sigma_0 - H_0) - (\Gamma - \varepsilon):
\]
the initial gap in stock, less the comparator's population discrepancy net of the UK's own
timing term. It is allocated between tenures by the comparator's private share of gross
building, $s = J_B^{\priv} / J_B$:
\begin{align*}
  \text{Private surplus} &= \bigl(J_N^{\priv} - N^{\priv}\bigr) + s\,\Delta,\\
  \text{Public surplus}  &= \bigl(J_N^{\pub} - N^{\pub}\bigr) + (1 - s)\,\Delta.
\end{align*}
The two add up to the total surplus. Country surpluses are measured against the UK's net
building $N^{\priv}$ and $N^{\pub}$ before the scaling in Part~I, Stage~4; the UK's timing
term $\varepsilon$ enters through $\Delta$ instead.

For the Netherlands, $s = 0.58968$ and
\[
  \Delta = (11{,}911.3 - 15{,}418.0) - (3{,}785.4 - 127.9) = -3{,}506.7 - 3{,}657.5 = -7{,}164.2,
\]
so
\begin{align*}
  \text{Private surplus} &= (13{,}248.6 - 8{,}137.7) + 0.58968 \times (-7{,}164.2)
                          = 5{,}110.9 - 4{,}224.6 = 886.3,\\
  \text{Public surplus}  &= (9{,}112.2 - 4{,}223.2) + 0.41032 \times (-7{,}164.2)
                          = 4{,}889.0 - 2{,}939.6 = 1{,}949.4.
\end{align*}

The results table reports these to four significant figures: for the Netherlands, 886,300
private, 1,949,000 public and 2,836,000 in total, with a private:public mix of gross
building of 59:41. The United Kingdom row reports $A^{\priv}$, $A^{\pub}$ and $\Delta H$
from Part~I (8,054,000, 4,179,000 and 12,230,000, a mix of 66:34, from
$N^{\priv}/N$). The Western European Average row is the simple mean of the eleven
comparators: 5,301,000 private, $-1{,}047{,}000$ public and 4,254,000 in total.

\subsection*{Notes}

\textbf{No floor on public additions.} The code applies no floor, and none binds: the
counterfactual's public net additions, $J_N^{\pub} + (1 - s)\Delta$, are positive for
every comparator. The smallest is Switzerland's, at 401,600.

\textbf{The published Table 3.} The report's Table 3 has the same total on every row as
this calculation but a different private/public split. Its UK row, 7,875,000 private and
4,358,000 public, splits total UK demolitions by the cumulative private share of gross
building instead of year by year:
\[
  G^{\priv} - (G - \Delta H)\,\frac{G^{\priv}}{G} = 7{,}875.
\]
That calculation is in the 2022 working code but not in
\code{Counterfactual Replication.R}. The published country rows, in which Switzerland and
Belgium show public surpluses of exactly $-4{,}358{,}000$, have not been reconciled with
any surviving code.

\textbf{Rounding.} Any difference between the worked figures here and the results table is
due to display rounding; the code rounds only its final values.

% ==============================================================================
\section*{Interpolation}

In the United Nations data, housing stock values are not given consistently across
Western Europe before 1960, and there are gaps of a year or more for certain countries.
The estimated stock $S_t$ fills them (\code{code/Summary Replication.R}, which reproduces
the stored series). Within each country, in date order:
\begin{enumerate}
  \item completions and demolitions are interpolated linearly across gaps of up to four
        years;
  \item the ratio of demolitions to completions is taken in every year that has both, and
        carried forward from the most recent such year, or back from the first, into years
        without;
  \item net additions are completions multiplied by one minus that ratio, and a year's
        stock includes that year's net additions.
\end{enumerate}
A country that never reports demolitions has no ratio, and its net additions are treated
as zero. Norway is the only such country among those in the report, so its estimate is a
straight line between reported stocks.

When there is a gap \textbf{between} years in reported housing stocks, the stock is rolled
forward from the earlier report on net additions, and the difference between the
rolled-forward value and the next report is spread evenly across the gap (linear
interpolation, for gaps of up to 35 years). For example, Austria reports 2.14 million
homes in 1951 and 2.37 million in 1963. Completions over 1952 to 1963 total 504,800.
Austria reports demolitions only from 1961, so the 1961 ratio of 2.5 per cent is carried
back to earlier years, giving 14,300 demolitions and 490,500 net additions. The reported
stock rose by only 229,600, so the shortfall of 260,900 is spread evenly: 21,700 a year
over the twelve years.

When there is no data \textbf{before} a given year's reported housing stock value, the
stock is rolled backwards from the first report on net additions. This matters for 1955
in two countries: Germany, whose first reported stock is for 1958, and Belgium, whose
first usable report is for 1963 (its 1948 figure is excluded as inconsistent with the
later series). Belgium reports 3.24 million homes in 1963. In 1962, 45,800 homes were
built and 2,000 demolished, a ratio of 4.4 per cent, which is carried back to earlier
years. Between 1956 and 1963, 365,000 homes were built and 18,600 demolished, leaving
346,400 net additions, which implies a housing stock of 2.89 million in 1955.

For historic England and Wales estimates, the demolitions between two census counts are
the difference between gross building and the change in the stock, distributed
\textbf{proportionally} to the gross building in each year. For example, the stock is
3,924,000 in 1861 and 4,520,000 in 1871, a rise of 596,000. As 749,000 homes were built in
1862 to 1871, this implies 153,000 demolitions. These are distributed over 1861 to 1870,
when 702,000 homes were built, so each year's demolitions are
\[
  \frac{\text{Annual gross building}}{702{,}000} \times 153{,}000.
\]

For graphs, methodological changes mean that discontinuities can arise -- for instance,
some countries' data appears to alternate between including and dropping vacant
dwellings. For example, after the UN data source ends in 2000, many countries switch to
counting inhabited dwellings only. In certain cases, the statistical authority will provide
inhabited data before 2000 as well, and this data, although not being used in the
calculations, is used in the graphs to avoid a discontinuity. If this is not possible, then
a consistent interpolated estimate of the housing stock using the first method is used for
graphs. This smoothed series is stored separately (\code{X5 (Alternative)}); the
calculations in this document use $S_t$ throughout.

For tenure, the public and private shares of construction are interpolated linearly across
gaps of up to five years. Beyond that, the most recent observed share is carried forward,
and before the first observation the first is carried back. The recorded value for
construction is then multiplied by these shares. In practice, large gaps in the tenure
data mainly occur for countries that build minimal social housing.

% ==============================================================================
\section*{Changes from the February 2023 annex}

\begin{itemize}[leftmargin=*]
  \item \textbf{Homes per person} divides by population; the published formula divided by
        the change in population.
  \item \textbf{Timing.} Building and demolitions are measured against the end-of-year
        stock, as in the code. This creates the term $\varepsilon$ in Part~I, which the
        code removes by scaling the UK's tenure totals, and a timing component in the
        Stage~6 discrepancy.
  \item \textbf{UK worked example.} The building rates are 2.13 and 1.98 per cent; the
        published 1.8 and 1.6 per cent were net of demolitions. Gross building,
        demolitions and net building by tenure are recomputed, and 58.6 per cent of
        demolitions, not two thirds, fall on private-built homes. UK tenure additions are
        8.05 and 4.18 million after scaling.
  \item \textbf{Population growth} is measured relative to the current year's population.
        The Netherlands' 1955 ratio is $1/1.008$, not $1/1.02$.
  \item \textbf{Demolition adjustment.} $D_t = d_t + (\xi_t / H_t - 1)/100$. The published
        formula had homes per person where a stock belongs and omitted the division by
        100. Negative values for Ireland are noted.
  \item \textbf{Stages 5--7} follow the code: gross building on the end-of-year
        counterfactual stock; the discrepancy defined as $J_N - (\nu_T - \sigma_0)$; and
        the remainder $\Delta$ allocated by the comparator's private share of gross
        building, against the UK's unscaled net building. The Netherlands example is
        recomputed. Its published gross building (23 million, of which 13.6 million
        private) matches the first-round stock $\xi_t$ rather than $\nu_t$.
  \item \textbf{Floor.} The published assumption that public additions stay positive is
        not applied by the code and does not bind. The note on the published Table 3 is new.
  \item \textbf{Interpolation.} Missing demolitions are filled with a carried ratio rather
        than by using gross completions; the Austria and Belgium examples are recomputed;
        Germany, as well as Belgium, is rolled back to 1955; Norway is a straight line; the
        England and Wales example uses the recorded 702,000 and 153,000; tenure gaps of up
        to five years, not four, are interpolated.
  \item \textbf{Tenure shares} that do not sum to one, for Finland and Switzerland, are
        noted.
\end{itemize}

% ==============================================================================
\section*{Acknowledgements}

Special thanks to Ronan Lyons, Mika Ronkainen and Per Spolander, who generously provided us
with post-2000 data for Ireland, Finland and Sweden. Special thanks also to Nick Crafts,
Paul Cheshire and John Myers for their advice early in the paper's research process.

\end{document}
